\section{Related literature}
\label{sec:literature}

\subsection{The economics of land value taxation}

The case for taxing land rests on the classical observation that land is supplied inelastically. \citet{george1879}, in \emph{Progress and Poverty}, argued that the rental value of unimproved land is a socially created surplus --- driven by population, public infrastructure and agglomeration rather than by any effort of the owner --- and proposed a ``single tax'' on land values to replace taxes on labour and capital. The modern optimal-taxation literature formalises the intuition: because the base cannot shrink in response to the tax, a levy on pure land rents generates no substitution effect and hence no deadweight loss, making it the canonical example of efficient taxation in the \citet{ramsey1927} inverse-elasticity framework. \citet{feldstein1977} qualified the classical view by noting general-equilibrium portfolio effects, but the presumption that land taxes are unusually efficient survives.

Two further strands reinforce the theoretical case. \citet{arnott1979} establish the Henry George Theorem: under conditions of optimal city size, aggregate land rents exactly finance the efficient level of local public expenditure, giving land taxation a normative role beyond efficiency --- it is the natural fiscal counterpart of the public goods that generate land value. In cross-country empirical work, \citet{johansson2008} rank recurrent taxes on immovable property as the least harmful to long-run growth among the major tax instruments, ahead of consumption, personal income and corporate income taxes. More recently, \citet{bonnet2021} document that the post-1980 rise in wealth-to-income ratios is to a large extent a rise in the price of residential land, and argue that land is accordingly both an increasingly important and an administratively feasible tax base for contemporary economies.

The empirical evidence, while limited by the scarcity of implemented land taxes, is broadly consistent with the theory. \citet{oates1997} examine Pittsburgh's two-rate property tax --- which taxed land at more than five times the rate applied to structures --- and find that the city sustained substantially higher building activity than comparable cities in the same period; they attribute this partly to the land tax's neutrality with respect to development, while noting the difficulty of isolating the tax's contribution from concurrent local conditions. \citet{tideman1995} argues on theoretical grounds that land taxation is better than neutral, as it discourages speculative withholding of developable land. The contributions collected in \citet{dye2009} temper this optimism by documenting the practical assessment challenges --- above all, separating land from improvement value at scale --- that have historically limited adoption.

\subsection{UK property taxation and council tax}

The deficiencies of council tax are well documented. \citet{mirrlees2011} concluded that UK housing taxation lacks a coherent rationale: council tax is regressive in property value and based on obsolete 1991 valuations, while stamp duty land tax taxes transactions and thereby impedes mobility. The review proposed replacing both with a Housing Services Tax proportional to the consumption value of housing, complemented by a land value tax on business land --- an explicit endorsement of land taxation where valuation is feasible. Subsequent work has quantified the regressivity that motivates reform: \citet{adam2020} show that effective council tax rates as a share of property value fall sharply across the price distribution, and that the absence of revaluation in England and Scotland leaves relative price changes since 1991 --- most notably London's appreciation --- untaxed. \citet{muellbauer2005} argues that reforming council tax into a proportional property tax with regular automatic revaluation would also stabilise the housing market, and the OECD's most recent survey of the UK likewise recommends updating property valuations and reforming council tax \citep{oecd2024}.

A body of applied UK work has modelled alternatives. \citet{corlett2018} examine options including a proportional property tax and a partial LVT at the aggregate and regional level; \citet{fairershare2021} propose a proportional property tax of 0.48 per cent of current property value and estimate that roughly three-quarters of households would pay less --- a useful comparator for our finding that 68 per cent of households gain at a 0.77 per cent tax on the narrower land base. \citet{ippr2019} model council tax reform for London. In Scotland, the cross-party Commission on Local Tax Reform \citep{scotgov2015} concluded that the present system must end and assessed property, land and income bases, and \citet{hughes2018} examine LVT options for the Scottish Land Commission, highlighting valuation and transition as the binding constraints. \citet{leunig2024} propose replacing both council tax and stamp duty with proportional levies on current values. Closest to our exercise, \citet{apgwilym2020} assess the feasibility of a local land value tax for Wales for the Welsh Government, simulating revenue-neutral rates at local-authority level using small-area land value estimates.

\subsection{Contribution}

Relative to this literature, our contribution is a household-level microsimulation of a full council-tax-to-LVT swap for the UK, with land values imputed to every household in a nationally representative micro-dataset, benchmarked against national balance-sheet aggregates, and evaluated within a complete model of the tax and benefit system. Prior UK LVT analyses have generally worked with regional, band-level or local-authority aggregates --- the Welsh assessment of \citet{apgwilym2020} being the most detailed to date --- rather than with household microdata. The microdata buy us three things that aggregate analyses cannot deliver: the joint distribution of land and income, which drives the winners-and-losers results and the divergence between income-ranked and wealth-ranked incidence; poverty rates measured against the statutory equivalised, housing-cost-adjusted thresholds; and a decomposition of who bears the corporate component of the base. We note in Section~\ref{sec:results} that the swap turns out to generate no interactions with means-tested benefit entitlements, so the household-level net income change is exactly the difference between the abolished council tax bill and the new land tax bill --- a finding, not an assumption, and one that materially simplifies interpretation. To our knowledge this is the first open-source, household-level distributional analysis of a revenue-neutral LVT for the whole of the UK, and the code and data pipeline are fully reproducible \citep{policyengine}.
